% Options for packages loaded elsewhere
\PassOptionsToPackage{unicode}{hyperref}
\PassOptionsToPackage{hyphens}{url}
\PassOptionsToPackage{dvipsnames,svgnames,x11names}{xcolor}
\documentclass[
  11pt,
]{article}
\usepackage{xcolor}
\usepackage[margin=1in]{geometry}
\usepackage{amsmath,amssymb}
\setcounter{secnumdepth}{-\maxdimen} % remove section numbering
\usepackage{iftex}
\ifPDFTeX
  \usepackage[T1]{fontenc}
  \usepackage[utf8]{inputenc}
  \usepackage{textcomp} % provide euro and other symbols
\else % if luatex or xetex
  \usepackage{unicode-math} % this also loads fontspec
  \defaultfontfeatures{Scale=MatchLowercase}
  \defaultfontfeatures[\rmfamily]{Ligatures=TeX,Scale=1}
\fi
\usepackage{lmodern}
\ifPDFTeX\else
  % xetex/luatex font selection
  \setmainfont[]{Georgia}
  \setmonofont[]{Menlo}
\fi
% Use upquote if available, for straight quotes in verbatim environments
\IfFileExists{upquote.sty}{\usepackage{upquote}}{}
\IfFileExists{microtype.sty}{% use microtype if available
  \usepackage[]{microtype}
  \UseMicrotypeSet[protrusion]{basicmath} % disable protrusion for tt fonts
}{}
\usepackage{setspace}
\makeatletter
\@ifundefined{KOMAClassName}{% if non-KOMA class
  \IfFileExists{parskip.sty}{%
    \usepackage{parskip}
  }{% else
    \setlength{\parindent}{0pt}
    \setlength{\parskip}{6pt plus 2pt minus 1pt}}
}{% if KOMA class
  \KOMAoptions{parskip=half}}
\makeatother
\usepackage{longtable,booktabs,array}
\newcounter{none} % for unnumbered tables
\usepackage{calc} % for calculating minipage widths
% Correct order of tables after \paragraph or \subparagraph
\usepackage{etoolbox}
\makeatletter
\patchcmd\longtable{\par}{\if@noskipsec\mbox{}\fi\par}{}{}
\makeatother
% Allow footnotes in longtable head/foot
\IfFileExists{footnotehyper.sty}{\usepackage{footnotehyper}}{\usepackage{footnote}}
\makesavenoteenv{longtable}
\usepackage{graphicx}
\makeatletter
\newsavebox\pandoc@box
\newcommand*\pandocbounded[1]{% scales image to fit in text height/width
  \sbox\pandoc@box{#1}%
  \Gscale@div\@tempa{\textheight}{\dimexpr\ht\pandoc@box+\dp\pandoc@box\relax}%
  \Gscale@div\@tempb{\linewidth}{\wd\pandoc@box}%
  \ifdim\@tempb\p@<\@tempa\p@\let\@tempa\@tempb\fi% select the smaller of both
  \ifdim\@tempa\p@<\p@\scalebox{\@tempa}{\usebox\pandoc@box}%
  \else\usebox{\pandoc@box}%
  \fi%
}
% Set default figure placement to htbp
\def\fps@figure{htbp}
\makeatother
\setlength{\emergencystretch}{3em} % prevent overfull lines
\providecommand{\tightlist}{%
  \setlength{\itemsep}{0pt}\setlength{\parskip}{0pt}}
\usepackage{bookmark}
\IfFileExists{xurl.sty}{\usepackage{xurl}}{} % add URL line breaks if available
\urlstyle{same}
\hypersetup{
  pdftitle={First-Pass Flood Binary Detection from Crowdsourced Imagery: Quality Control for Screening Systems},
  colorlinks=true,
  linkcolor={Maroon},
  filecolor={Maroon},
  citecolor={Blue},
  urlcolor={Blue},
  pdfcreator={LaTeX via pandoc}}

\title{First-Pass Flood Binary Detection from Crowdsourced Imagery:
Quality Control for Screening Systems}
\author{}
\date{2026-04-22}

\begin{document}
\maketitle

\setstretch{1.35}
\begin{quote}
\textbf{One-sentence contribution:} We show that Phase-1 Hard Negative
Mining --- mining confounder candidates from a partially-trained
intermediate checkpoint rather than the converged model ---
simultaneously improves flood recall (97.8\%→99.1\%, FN 7→3) and reduces
the dominant river false positive rate (9.2\%→5.26\%) for
EfficientNetB0-based flood screening on crowdsourced street-level
imagery, outperforming extended training without data augmentation.
\end{quote}

\begin{center}\rule{0.5\linewidth}{0.5pt}\end{center}

\section{Abstract}\label{abstract}

We introduce a \textbf{Phase-1 Hard Negative Mining (HNM)} protocol for
improving flood/non-flood binary classifiers trained on crowdsourced
street-level imagery. Automated flood screening systems routinely
false-positive on visually confounding non-flood scenes --- swimming
pools, rivers, and wet roads share water texture, reflectance, and urban
context with flooded streets --- yet no prior work applies targeted
confounder resampling to this problem, and existing evaluations report
aggregate precision rather than per-category false positive rates. We
train two-phase progressively fine-tuned EfficientNetB0 and ResNet50
baselines under both binary cross-entropy and focal loss, then mine hard
negatives specifically from the Phase 1 checkpoint --- where the
backbone is still largely frozen and confounder confusion is highest ---
and inject them into retraining from the fully converged Phase 2
checkpoint. We evaluate on 4,099 crowdsourced images spanning 8
confounder categories, reporting PR-AUC as the primary metric alongside
per-category false positive rates with Clopper-Pearson confidence
intervals. Both architectures achieve strong baseline performance:
EfficientNetB0 BCE reaches PR-AUC \textbf{0.9976}, accuracy
\textbf{98.3\%}, and flood recall \textbf{97.8\%} (7 FN, 1 in 46), while
ResNet50 BCE reaches PR-AUC \textbf{0.9991}, accuracy \textbf{98.5\%},
and flood recall \textbf{97.2\%} (9 FN, 1 in 36). Despite ResNet50's
marginally higher PR-AUC, EfficientNetB0 achieves better per-threshold
flood recall while incurring a higher river false positive rate (9.2\%
vs.~3.9\%), creating the primary target for HNM. Phase-1 HNM applied to
EfficientNetB0 BCE achieves \textbf{99.1\%} flood recall with only
\textbf{3 missed floods} (1 in 107), best overall accuracy
(\textbf{98.78\%}), and reduces river false positives from \textbf{9.2\%
to 5.26\%} (7→4 FP), combining the recall advantage of EfficientNetB0
with a river FP rate approaching ResNet50's baseline.

\begin{center}\rule{0.5\linewidth}{0.5pt}\end{center}

\section{1. Introduction}\label{introduction}

Floods are among the most destructive and frequently occurring natural
disasters worldwide, causing significant loss of life and infrastructure
damage. Rapid automated detection of flooding from crowdsourced imagery
--- photographs submitted via smartphone apps, social media, or traffic
camera feeds --- can dramatically accelerate emergency response, often
preceding official reports by hours (Esparza et al., 2022). The
first-pass screening problem is binary: given an image from an unknown
contributor, does it depict active flooding?

\textbf{The screening objective is fundamentally recall-first.} A missed
flood image (false negative) means an affected location is not escalated
for human review; a false positive triggers unnecessary resource
allocation. The asymmetric cost places flood recall above precision as
the primary operational objective. Yet most published flood detection
systems report ROC-AUC or overall accuracy, which are known to overstate
performance on imbalanced binary classification (Davis \& Goadrich,
2006) and do not reflect the recall-constrained operating point required
for deployment.

\textbf{Visual confounders are the dominant failure mode.} Swimming
pools, rivers, and wet roads share an open reflective water surface,
similar colour statistics under overcast conditions, and proximity to
urban infrastructure --- precisely the features a texture-based
convolutional classifier uses to detect floods. Under close-range
photography, contextual cues that would disambiguate a pool from a flood
are absent, and standard cross-entropy trained models frequently
misclassify these scenes. This problem is underexplored: existing flood
detection datasets rarely include systematic confounder coverage, and
few papers report per-category false positive rates separately from
aggregate precision.

\textbf{Research gap.} Flood detection systems achieve high aggregate
accuracy --- Lee et al.~(2025) report 97.67\% recall; Khan et al.~(2023)
report 96.43\% accuracy --- yet none measure or address per-category
confounder false positive rates. No published flood screening work (a)
evaluates FP rates broken down by confounder category, (b) applies hard
negative mining from an intermediate checkpoint, or (c) uses the Phase
1/Phase 2 distinction to recover meaningful mining signal after
convergence. This matters operationally: a system with 97\% flood recall
but 30\% FP rate on rivers will generate massive alert fatigue in
flooded urban areas --- precisely where rivers and floods co-occur.

\textbf{We address both gaps.} First, we build a benchmark with explicit
confounder category labels and measure per-category false positive rates
on a held-out validation set. Second, we implement \textbf{offline Hard
Negative Mining}: we identify the hardest-scoring non-flood images using
the \emph{Phase 1} checkpoint --- before the model has memorised the
training distribution --- and inject them into retraining from the fully
converged Phase 2 baseline. This Phase 1 mining design is a novel
contribution: by end of Phase 2, FP rates on training data collapse to
near zero, leaving no candidates; the Phase 1 checkpoint retains genuine
confusion on visually ambiguous categories.

\textbf{Contributions:}

\begin{enumerate}
\def\labelenumi{\arabic{enumi}.}
\tightlist
\item
  A flood/non-flood benchmark with fine-grained labels across 8
  confounder categories with a clean stratified 80/20 train/val split
  (4,099 images after deduplication via SHA-256 hashing).
\item
  A two-phase progressive fine-tuning protocol (EfficientNetB0 and
  ResNet50, BCE and focal loss) with class-weighted training,
  establishing that both architectures achieve competitive performance
  (PR-AUC ≥ 0.997) when preprocessing is correctly configured.
\item
  \textbf{Phase-1 HNM}: confounder analysis from the Phase 1 checkpoint,
  not the converged Phase 2 model, producing genuine mining candidates
  for augmented retraining --- a design that avoids the memorisation
  problem that defeats Phase-2 mining.
\item
  Ablation controls: no-injection (extended training without HNM) and
  random-injection (same count, probability-unranked), isolating the
  contribution of difficulty-based candidate selection.
\item
  Per-category false positive rate analysis (8 categories,
  Clopper-Pearson 95\% CI), identifying river as the sole significant
  confounder across all conditions and quantifying the statistical power
  required for meaningful inference.
\end{enumerate}

\begin{center}\rule{0.5\linewidth}{0.5pt}\end{center}

\section{2. Related Work}\label{related-work}

{[}PLACEHOLDER --- insert from existing paper{]}

\begin{center}\rule{0.5\linewidth}{0.5pt}\end{center}

\section{3. Dataset}\label{dataset}

\textbf{Sources.} We combine two public datasets:

\begin{itemize}
\tightlist
\item
  \textbf{USF FloodingDataset2}: 1,613 flood images labeled by severity
  (Major, Moderate, Minor) and 2,087 non-flood images spanning 7 visual
  confounder categories (river, swimming pool, wet/rain-slicked road,
  park/walkway, building, vehicle, animal, plant).
\item
  \textbf{RIWA (Wagner et al., 2023)}: 399 additional river images,
  substantially enriching the dominant confounder category from its
  natural representation in USF FloodingDataset2.
\end{itemize}

\textbf{Deduplication.} SHA-256 hashing identifies and removes 55
exact-duplicate images from the initial 4,154-image corpus, yielding
\textbf{4,099 unique images}.

\textbf{Split.} We apply stratified 80/20 splitting (seed 42) that
preserves both the overall flood prevalence (39.3\%) and the
per-category distribution within the non-flood class:

{\def\LTcaptype{none} % do not increment counter
\begin{longtable}[]{@{}lllll@{}}
\toprule\noalign{}
Split & Total & Flood & Non-Flood & Flood Prevalence \\
\midrule\noalign{}
\endhead
\bottomrule\noalign{}
\endlastfoot
Train & 3,280 & 1,290 & 1,990 & 39.3\% \\
Val & 819 & 322 & 497 & 39.3\% \\
\end{longtable}
}

\textbf{Confounder categories (validation set):}

{\def\LTcaptype{none} % do not increment counter
\begin{longtable}[]{@{}
  >{\raggedright\arraybackslash}p{(\linewidth - 4\tabcolsep) * \real{0.3333}}
  >{\raggedright\arraybackslash}p{(\linewidth - 4\tabcolsep) * \real{0.3333}}
  >{\raggedright\arraybackslash}p{(\linewidth - 4\tabcolsep) * \real{0.3333}}@{}}
\toprule\noalign{}
\begin{minipage}[b]{\linewidth}\raggedright
Category
\end{minipage} & \begin{minipage}[b]{\linewidth}\raggedright
N Val Images
\end{minipage} & \begin{minipage}[b]{\linewidth}\raggedright
Notes
\end{minipage} \\
\midrule\noalign{}
\endhead
\bottomrule\noalign{}
\endlastfoot
river & 76 & Primary confounder; enriched via RIWA \\
street\_clear & 105 & Largest non-flood category \\
park\_walkway & 82 & --- \\
animal & 67 & --- \\
building & 55 & --- \\
vehicle & 50 & --- \\
plant & 34 & --- \\
swimming\_pool & 28 & Statistically underpowered; upper 95\% CI = 12.3\%
at 0\% observed FP \\
\end{longtable}
}

Val images are never included in mining candidates. A partition-safety
assertion is run before each HNM retraining to verify there is zero
train/val overlap.

\begin{center}\rule{0.5\linewidth}{0.5pt}\end{center}

\section{4. Methods}\label{methods}

\subsection{4.1 Two-Phase Progressive
Fine-Tuning}\label{two-phase-progressive-fine-tuning}

All models use \textbf{two-phase transfer learning} from
ImageNet-pretrained weights. Progressive fine-tuning mitigates
catastrophic forgetting during fine-tuning (Kirkpatrick et al., 2017;
Lyu et al., 2025):

\begin{itemize}
\tightlist
\item
  \textbf{Phase 1} --- Freeze all but the last \texttt{n\_trainable=30}
  backbone layers. Train classification head and top backbone layers at
  LR=1×10⁻³, max 30 epochs.
\item
  \textbf{Phase 2} --- Freeze the first \texttt{n\_frozen=50} backbone
  layers; fine-tune deeper layers at LR=1×10⁻⁴, max 50 epochs.
\end{itemize}

Both phases use early stopping (patience=10 on val loss) and
ReduceLROnPlateau (patience=5 Phase 1, patience=3 Phase 2). Class
weights are computed per run via
\texttt{compute\_class\_weight("balanced")} on the training label
distribution. The model is saved at the best val loss epoch in each
phase.

\begin{quote}
\textbf{Note on training recall metric}: Keras's \texttt{Recall} metric
during training measures recall for class index 1 (non\_flood,
alphabetically assigned by \texttt{flow\_from\_directory}). Training
curves therefore report \textbf{non-flood} recall during training. Flood
recall is computed from the confusion matrix in post-training
evaluation.
\end{quote}

\subsection{4.2 Architectures}\label{architectures}

{\def\LTcaptype{none} % do not increment counter
\begin{longtable}[]{@{}
  >{\raggedright\arraybackslash}p{(\linewidth - 6\tabcolsep) * \real{0.2500}}
  >{\raggedright\arraybackslash}p{(\linewidth - 6\tabcolsep) * \real{0.2500}}
  >{\raggedright\arraybackslash}p{(\linewidth - 6\tabcolsep) * \real{0.2500}}
  >{\raggedright\arraybackslash}p{(\linewidth - 6\tabcolsep) * \real{0.2500}}@{}}
\toprule\noalign{}
\begin{minipage}[b]{\linewidth}\raggedright
Model
\end{minipage} & \begin{minipage}[b]{\linewidth}\raggedright
Params
\end{minipage} & \begin{minipage}[b]{\linewidth}\raggedright
Input preprocessing
\end{minipage} & \begin{minipage}[b]{\linewidth}\raggedright
Notes
\end{minipage} \\
\midrule\noalign{}
\endhead
\bottomrule\noalign{}
\endlastfoot
EfficientNetB0 & \textasciitilde5.3M & Raw {[}0, 255{]} --- built-in
rescaling layer & Keras 3: \texttt{include\_preprocessing} removed;
preprocessing is in the model graph \\
ResNet50 & \textasciitilde25.6M & Channel mean subtraction (ImageNet
stats) & Standard
\texttt{tf.keras.applications.resnet50.preprocess\_input} \\
\end{longtable}
}

The preprocessing difference is critical and verified at runtime.
EfficientNet receives raw {[}0, 255{]} pixel values processed by its
internal scaling layer; ResNet50 receives ImageNet-mean-subtracted
inputs. Mismatched preprocessing causes dramatic accuracy degradation
and should be treated as a model-architecture pairing constraint, not a
tunable hyperparameter.

\subsection{4.3 Loss Functions}\label{loss-functions}

\begin{itemize}
\tightlist
\item
  \textbf{Binary Cross-Entropy (BCE)}: Standard sigmoid loss. Class
  weights compensate for global label imbalance (flood:non-flood ≈ 1:1.5
  in train set).
\item
  \textbf{Focal Loss} (γ=2.0): (1 − p\_t)\^{}γ × BCE. Down-weights
  well-classified examples within each mini-batch, focusing gradient
  updates on hard cases. Also trained with balanced class weights.
\end{itemize}

Focal loss and class weighting address different aspects of the learning
problem: class weighting corrects for global class frequency imbalance;
focal loss corrects for local within-batch hard/easy example imbalance.

\subsection{4.4 Phase-1 Hard Negative Mining (Core
Contribution)}\label{phase-1-hard-negative-mining-core-contribution}

The HNM pipeline proceeds in three stages:

\subsubsection{Stage 1 --- Confounder analysis from the Phase 1
checkpoint}\label{stage-1-confounder-analysis-from-the-phase-1-checkpoint}

Run the \textbf{Phase 1} checkpoint on all \texttt{train/non\_flood/}
images. For each image, compute flood probability p̂ = 1 − σ(output),
where the sigmoid output is P(non\_flood) due to alphabetical class
ordering in \texttt{flow\_from\_directory}. Compute per-category FP rate
at τ=0.5; flag categories exceeding θ=0.05.

\textbf{Why Phase 1?} By end of Phase 2, the model fits the training
non-flood categories to near-zero FP rate on the training set --- there
are no binary FPs to flag and no candidates to mine. The observed Phase
2 training-set river FP rate is 0.31\% for EfficientNetB0 ---
essentially no signal. The Phase 1 checkpoint has only trained the
classification head and the top 30 backbone layers; the backbone is
still largely frozen, leaving residual confusion on visually ambiguous
categories. Mining from Phase 1 produces genuine candidates while Phase
2 remains the strong starting point for HNM retraining. This is
analogous to knowledge distillation, where intermediate representations
often retain more generalisable uncertainty than the converged model.

Partition safety: val images are never included in mining candidates. A
leakage assertion is run before retraining begins.

\subsubsection{Stage 2 --- Candidate selection by flood
probability}\label{stage-2-candidate-selection-by-flood-probability}

From flagged categories, collect all images. Rank by descending Phase 1
flood probability p̂. Select the \textbf{top 20\%} (percentile mode) ---
the images the model is most confused about at the Phase 1 stage, not
just those that crossed the binary threshold.

\subsubsection{Stage 3 --- Augmented retraining from the Phase 2
checkpoint}\label{stage-3-augmented-retraining-from-the-phase-2-checkpoint}

Inject selected hard negatives into the training set. Apply 5×
augmentation to each selected image. Retrain from the \textbf{Phase 2}
checkpoint at Phase 2 LR (1×10⁻⁴). The Phase 2 checkpoint is the strong
converged baseline; HNM retraining refines its decision boundary
specifically in the confounder regions.

\subsubsection{Ablation Controls}\label{ablation-controls}

{\def\LTcaptype{none} % do not increment counter
\begin{longtable}[]{@{}
  >{\raggedright\arraybackslash}p{(\linewidth - 4\tabcolsep) * \real{0.3333}}
  >{\raggedright\arraybackslash}p{(\linewidth - 4\tabcolsep) * \real{0.3333}}
  >{\raggedright\arraybackslash}p{(\linewidth - 4\tabcolsep) * \real{0.3333}}@{}}
\toprule\noalign{}
\begin{minipage}[b]{\linewidth}\raggedright
Condition
\end{minipage} & \begin{minipage}[b]{\linewidth}\raggedright
Description
\end{minipage} & \begin{minipage}[b]{\linewidth}\raggedright
Controls for
\end{minipage} \\
\midrule\noalign{}
\endhead
\bottomrule\noalign{}
\endlastfoot
\textbf{HNM (percentile)} & Top 10\% hardest negatives by Phase 1
probability & --- \\
\textbf{No injection} & Retrain for same epoch budget, no additional
data & Additional training epochs without distribution change \\
\textbf{Random injection} & Same count, randomly sampled from flagged
categories & Category exposure without difficulty ranking \\
\end{longtable}
}

The no-injection control is critical: if HNM improves performance, it
must outperform extended training without injection to attribute the
gain to difficulty-ranked hard negative selection rather than additional
compute.

\subsection{4.5 Metrics and Evaluation
Protocol}\label{metrics-and-evaluation-protocol}

\textbf{Primary metric: PR-AUC} (area under the precision-recall curve).
PR-AUC is preferred over ROC-AUC for imbalanced screening problems
(Davis \& Goadrich, 2006; Saito \& Rehmsmeier, 2015): ROC-AUC inflates
performance because specificity is easy to achieve when negatives
substantially outnumber positives at the operating threshold.

\textbf{Threshold selection}: τ is selected on the val set to identify
the recall-preserving operating range. Unless otherwise noted, τ=0.5 is
used for per-threshold comparisons. Test-set threshold tuning is not
performed.

\textbf{Statistical analysis}: - Model comparisons: McNemar's test on
paired predictions (Bonferroni-corrected) - Per-category FP rates:
Clopper-Pearson 95\% CI (exact binomial); Fisher's exact test for
pre/post-HNM comparisons per category - Multi-seed bootstrap CIs pending
(5 seeds: 42, 1, 7, 123, 2024)

\subsection{4.6 Probability Calibration}\label{probability-calibration}

Modern deep classifiers are overconfident (Guo et al., 2017). We apply
\textbf{temperature scaling} post-hoc --- dividing logits by a scalar T
learned by minimising NLL on the val set. Temperature scaling introduces
one degree of freedom, making it the most stable calibration method on
small validation sets. Calibration quality is reported via Expected
Calibration Error (ECE) and reliability diagrams. \emph{{[}Calibration
results pending.{]}}

\begin{center}\rule{0.5\linewidth}{0.5pt}\end{center}

\section{5. Experiments}\label{experiments}

\subsection{5.1 Baseline Training
Dynamics}\label{baseline-training-dynamics}

Figure 1 shows training curves for the EfficientNetB0 BCE baseline. The
Phase 1/Phase 2 boundary is visible around epoch 7--8 as a spike in
training loss --- when the frozen backbone layers are unfrozen and the
learning rate is reduced, the model briefly increases training loss
before converging. Validation loss (dashed) remains stable throughout,
with a train-val gap of \textasciitilde0.02 indicating mild overfitting.
AUC and recall are high and stable across both phases.

\begin{figure}
\centering
\pandocbounded{\includegraphics[keepaspectratio,alt={EfficientNetB0 BCE training curves --- loss, AUC, accuracy, and recall per epoch.}]{../results/figures/baselines/efficientnet_bce/training_curves.png}}
\caption{EfficientNetB0 BCE training curves --- loss, AUC, accuracy, and
recall per epoch.}
\end{figure}

\emph{Figure 1a: EfficientNetB0 BCE baseline training curves (seed 42).
Dashed = validation; solid = training. Phase 1→Phase 2 transition
visible as training loss spike at epoch \textasciitilde7.}

\begin{figure}
\centering
\pandocbounded{\includegraphics[keepaspectratio,alt={ResNet50 BCE training curves --- loss, AUC, accuracy, and recall per epoch.}]{../results/figures/baselines/resnet50_bce/training_curves.png}}
\caption{ResNet50 BCE training curves --- loss, AUC, accuracy, and
recall per epoch.}
\end{figure}

\emph{Figure 1b: ResNet50 BCE baseline training curves (seed 42). Best
checkpoint at Phase 2 epoch 16 (⭐). Phase 2 converges in fewer epochs
than EfficientNetB0 due to early stopping triggered by val loss
plateau.}

\subsection{5.2 Baseline Results at τ =
0.5}\label{baseline-results-at-ux3c4-0.5}

Table 1 reports val set performance for all four baseline conditions
(seed 42). Both EfficientNetB0 and ResNet50 BCE achieve strong
performance when correctly configured.

{\def\LTcaptype{none} % do not increment counter
\begin{longtable}[]{@{}
  >{\raggedright\arraybackslash}p{(\linewidth - 14\tabcolsep) * \real{0.1250}}
  >{\raggedright\arraybackslash}p{(\linewidth - 14\tabcolsep) * \real{0.1250}}
  >{\raggedright\arraybackslash}p{(\linewidth - 14\tabcolsep) * \real{0.1250}}
  >{\raggedright\arraybackslash}p{(\linewidth - 14\tabcolsep) * \real{0.1250}}
  >{\raggedright\arraybackslash}p{(\linewidth - 14\tabcolsep) * \real{0.1250}}
  >{\raggedright\arraybackslash}p{(\linewidth - 14\tabcolsep) * \real{0.1250}}
  >{\raggedright\arraybackslash}p{(\linewidth - 14\tabcolsep) * \real{0.1250}}
  >{\raggedright\arraybackslash}p{(\linewidth - 14\tabcolsep) * \real{0.1250}}@{}}
\toprule\noalign{}
\begin{minipage}[b]{\linewidth}\raggedright
Model
\end{minipage} & \begin{minipage}[b]{\linewidth}\raggedright
Loss
\end{minipage} & \begin{minipage}[b]{\linewidth}\raggedright
Val Acc ↑
\end{minipage} & \begin{minipage}[b]{\linewidth}\raggedright
Flood Recall ↑
\end{minipage} & \begin{minipage}[b]{\linewidth}\raggedright
Flood Prec ↑
\end{minipage} & \begin{minipage}[b]{\linewidth}\raggedright
Flood F1 ↑
\end{minipage} & \begin{minipage}[b]{\linewidth}\raggedright
FN ↓
\end{minipage} & \begin{minipage}[b]{\linewidth}\raggedright
FP ↓
\end{minipage} \\
\midrule\noalign{}
\endhead
\bottomrule\noalign{}
\endlastfoot
EfficientNetB0 & BCE & 98.3\% & 97.8\% & 0.978 & 0.978 & 7 (1 in 46) &
7 \\
\textbf{EfficientNetB0} & \textbf{Focal} & \textbf{98.8\%} &
\textbf{98.5\%} & \textbf{0.990} & \textbf{0.985} & \textbf{5 (1 in 64)}
& \textbf{5} \\
ResNet50 & BCE & \textbf{98.5\%} & 97.2\% & 0.991 & 0.981 & 9 (1 in 36)
& 3 \\
ResNet50 & Focal & \emph{pending} & \emph{pending} & --- & --- & --- &
--- \\
\end{longtable}
}

\emph{Table 1: Baseline val set performance, seed 42, τ=0.5. Both
architectures are competitive when preprocessing is correctly matched to
architecture requirements. EfficientNetB0 achieves higher per-threshold
flood recall (7 vs.~9 FN) while ResNet50 achieves lower false positives
(3 vs.~7 FP). Focal loss improves EfficientNetB0 (97.8\%→98.5\% recall,
7→5 FN); ResNet50 Focal results pending re-evaluation.}

{\def\LTcaptype{none} % do not increment counter
\begin{longtable}[]{@{}
  >{\raggedright\arraybackslash}p{(\linewidth - 8\tabcolsep) * \real{0.2000}}
  >{\raggedright\arraybackslash}p{(\linewidth - 8\tabcolsep) * \real{0.2000}}
  >{\raggedright\arraybackslash}p{(\linewidth - 8\tabcolsep) * \real{0.2000}}
  >{\raggedright\arraybackslash}p{(\linewidth - 8\tabcolsep) * \real{0.2000}}
  >{\raggedright\arraybackslash}p{(\linewidth - 8\tabcolsep) * \real{0.2000}}@{}}
\toprule\noalign{}
\begin{minipage}[b]{\linewidth}\raggedright
Model
\end{minipage} & \begin{minipage}[b]{\linewidth}\raggedright
Loss
\end{minipage} & \begin{minipage}[b]{\linewidth}\raggedright
\textbf{PR-AUC ↑}
\end{minipage} & \begin{minipage}[b]{\linewidth}\raggedright
ROC-AUC ↑
\end{minipage} & \begin{minipage}[b]{\linewidth}\raggedright
Notes
\end{minipage} \\
\midrule\noalign{}
\endhead
\bottomrule\noalign{}
\endlastfoot
EfficientNetB0 & BCE & 0.9976 & 0.9985 & --- \\
\textbf{EfficientNetB0} & \textbf{Focal} & \textbf{0.9977} &
\textbf{0.9986} & Marginal PR-AUC gain; clear per-threshold recall
gain \\
ResNet50 & BCE & \textbf{0.9991} & \textbf{0.9995} & Higher PR-AUC;
fewer FP; but 2 more FN than EfficientNetB0 BCE \\
ResNet50 & Focal & \emph{pending} & \emph{pending} & --- \\
\end{longtable}
}

\emph{Table 1b: Discrimination metrics (post-hoc sklearn, val set, seed
42). ResNet50 BCE achieves marginally higher PR-AUC (0.9991 vs.~0.9976)
and ROC-AUC, while EfficientNetB0 BCE achieves better per-threshold
flood recall at τ=0.5 (7 vs.~9 FN). The two architectures present a
meaningful tradeoff: EfficientNetB0 prioritises recall (misses fewer
floods at τ=0.5) while ResNet50 prioritises precision on non-floods (3
vs.~7 FP). For recall-first screening, this tradeoff favours
EfficientNetB0 at the default operating point.}

The confusion matrices (Figure 2) make the tradeoff concrete.
EfficientNetB0 BCE: 7 FN, 7 FP --- symmetric errors. ResNet50 BCE: 9 FN,
3 FP --- asymmetric, biased toward predicting non-flood. EfficientNetB0
Focal achieves the best per-threshold result with only 5 FN and 5 FP.

\begin{figure}
\centering
\pandocbounded{\includegraphics[keepaspectratio,alt={EfficientNetB0 BCE confusion matrix (TP=315, FN=7, FP=7, TN=490)}]{../results/figures/baselines/efficientnet_bce/confusion_matrix.png}}
\caption{EfficientNetB0 BCE confusion matrix (TP=315, FN=7, FP=7,
TN=490)}
\end{figure}

\emph{Figure 2a: EfficientNetB0 BCE confusion matrix (val, n=819, seed
42). Near-symmetric errors: 7 FN, 7 FP.}

\begin{figure}
\centering
\pandocbounded{\includegraphics[keepaspectratio,alt={ResNet50 BCE confusion matrix (TP=313, FN=9, FP=3, TN=494)}]{../results/figures/baselines/resnet50_bce/confusion_matrix.png}}
\caption{ResNet50 BCE confusion matrix (TP=313, FN=9, FP=3, TN=494)}
\end{figure}

\emph{Figure 2b: ResNet50 BCE confusion matrix (val, n=819, seed 42).
Asymmetric: 9 FN, 3 FP. Prioritises precision on non-floods over flood
recall.}

\subsection{5.3 PR-AUC vs.~ROC-AUC}\label{pr-auc-vs.-roc-auc}

Figure 3 shows PR and ROC curves for both BCE baselines. At the curve
level, ResNet50 achieves slightly higher PR-AUC (0.9991 vs.~0.9976),
reflecting strong precision across all recall operating points.
EfficientNetB0 achieves better recall at the τ=0.5 operating point
(97.8\% vs.~97.2\%, 7 vs.~9 FN). For deployment at a fixed threshold,
EfficientNetB0 is the better choice; at a variable threshold calibrated
to a recall target, ResNet50's higher PR-AUC offers more precision
headroom. Both architectures dramatically outperform what a single
aggregate metric would suggest, reinforcing Davis \& Goadrich (2006)'s
argument that PR-AUC is the appropriate primary metric for this problem.

\begin{figure}
\centering
\pandocbounded{\includegraphics[keepaspectratio,alt={EfficientNetB0 PR and ROC curves --- PR AUC=0.9976, ROC AUC=0.9985}]{../results/figures/baselines/efficientnet_bce/pr_roc_curves.png}}
\caption{EfficientNetB0 PR and ROC curves --- PR AUC=0.9976, ROC
AUC=0.9985}
\end{figure}

\emph{Figure 3a: EfficientNetB0 BCE precision-recall and ROC curves (val
set, seed 42). PR-AUC=0.9976.}

\begin{figure}
\centering
\pandocbounded{\includegraphics[keepaspectratio,alt={ResNet50 PR and ROC curves --- PR AUC=0.9991, ROC AUC=0.9995}]{../results/figures/baselines/resnet50_bce/pr_roc_curves.png}}
\caption{ResNet50 PR and ROC curves --- PR AUC=0.9991, ROC AUC=0.9995}
\end{figure}

\emph{Figure 3b: ResNet50 BCE precision-recall and ROC curves (val set,
seed 42). PR-AUC=0.9991. Marginally higher than EfficientNetB0, but
achieves fewer FN only at the curve level; at τ=0.5, EfficientNetB0
recalls more floods.}

\subsection{5.4 Qualitative Error
Analysis}\label{qualitative-error-analysis}

Figure 4 shows the worst FP (top, red border) and FN (bottom, orange
border) examples for each model at τ=0.5. EfficientNetB0's errors are
genuinely ambiguous boundary cases --- river scenes with strong water
texture, and dimly-lit or partially obscured flood scenes. ResNet50's
fewer but different FN cases are harder to characterise at N=9;
qualitative review of the 9 missed floods in ResNet50's error set
reveals a mix of genuinely ambiguous and mildly clear scenes, consistent
with a model whose decision boundary is well-calibrated in aggregate but
slightly conservative at τ=0.5.

\begin{figure}
\centering
\pandocbounded{\includegraphics[keepaspectratio,alt={EfficientNetB0 worst FP and FN examples}]{../results/figures/baselines/efficientnet_bce/fp_fn_examples.png}}
\caption{EfficientNetB0 worst FP and FN examples}
\end{figure}

\emph{Figure 4a: EfficientNetB0 BCE worst val errors (seed 42). Top row
(red): false positives --- visually ambiguous river/water scenes. Bottom
row (orange): false negatives --- dimly lit or partially obscured flood
scenes.}

\begin{figure}
\centering
\pandocbounded{\includegraphics[keepaspectratio,alt={ResNet50 worst FP and FN examples}]{../results/figures/baselines/resnet50_bce/fp_fn_examples.png}}
\caption{ResNet50 worst FP and FN examples}
\end{figure}

\emph{Figure 4b: ResNet50 BCE worst val errors (seed 42). Top row (red):
3 false positives. Bottom row (orange): 9 false negatives --- a mix of
ambiguous and moderately clear flood scenes, consistent with a slightly
conservative decision boundary.}

\subsection{5.5 Confounder Analysis}\label{confounder-analysis}

Table 2 reports per-category FP rates on the val split for both BCE
baselines. \textbf{River is the sole significant confounder across both
architectures.} All other 7 categories observe 0\% FP with tight upper
confidence bounds.

{\def\LTcaptype{none} % do not increment counter
\begin{longtable}[]{@{}
  >{\raggedright\arraybackslash}p{(\linewidth - 14\tabcolsep) * \real{0.1250}}
  >{\raggedright\arraybackslash}p{(\linewidth - 14\tabcolsep) * \real{0.1250}}
  >{\raggedright\arraybackslash}p{(\linewidth - 14\tabcolsep) * \real{0.1250}}
  >{\raggedright\arraybackslash}p{(\linewidth - 14\tabcolsep) * \real{0.1250}}
  >{\raggedright\arraybackslash}p{(\linewidth - 14\tabcolsep) * \real{0.1250}}
  >{\raggedright\arraybackslash}p{(\linewidth - 14\tabcolsep) * \real{0.1250}}
  >{\raggedright\arraybackslash}p{(\linewidth - 14\tabcolsep) * \real{0.1250}}
  >{\raggedright\arraybackslash}p{(\linewidth - 14\tabcolsep) * \real{0.1250}}@{}}
\toprule\noalign{}
\begin{minipage}[b]{\linewidth}\raggedright
Category
\end{minipage} & \begin{minipage}[b]{\linewidth}\raggedright
N Val
\end{minipage} & \begin{minipage}[b]{\linewidth}\raggedright
EffNet FP
\end{minipage} & \begin{minipage}[b]{\linewidth}\raggedright
EffNet Rate
\end{minipage} & \begin{minipage}[b]{\linewidth}\raggedright
95\% CP CI
\end{minipage} & \begin{minipage}[b]{\linewidth}\raggedright
ResNet50 FP
\end{minipage} & \begin{minipage}[b]{\linewidth}\raggedright
ResNet50 Rate
\end{minipage} & \begin{minipage}[b]{\linewidth}\raggedright
95\% CP CI
\end{minipage} \\
\midrule\noalign{}
\endhead
\bottomrule\noalign{}
\endlastfoot
\textbf{river} & 76 & 7 & \textbf{9.21\%} & \textbf{{[}3.8\%, 17.7\%{]}}
& 3 & \textbf{3.95\%} & \textbf{{[}0.8\%, 11.0\%{]}} \\
plant & 34 & 1 & 2.94\% & {[}0.1\%, 15.3\%{]} & 0 & 0.0\% & {[}0.0\%,
10.3\%{]} \\
animal & 67 & 0 & 0.0\% & {[}0.0\%, 5.4\%{]} & 0 & 0.0\% & {[}0.0\%,
5.4\%{]} \\
building & 55 & 0 & 0.0\% & {[}0.0\%, 6.5\%{]} & 0 & 0.0\% & {[}0.0\%,
6.5\%{]} \\
park\_walkway & 82 & 0 & 0.0\% & {[}0.0\%, 4.4\%{]} & 0 & 0.0\% &
{[}0.0\%, 4.4\%{]} \\
street\_clear & 105 & 0 & 0.0\% & {[}0.0\%, 3.5\%{]} & 0 & 0.0\% &
{[}0.0\%, 3.5\%{]} \\
swimming\_pool & 28 & 0 & 0.0\% & {[}0.0\%, \textbf{12.3\%}{]} & 0 &
0.0\% & {[}0.0\%, \textbf{12.3\%}{]} \\
vehicle & 50 & 0 & 0.0\% & {[}0.0\%, 7.1\%{]} & 0 & 0.0\% & {[}0.0\%,
7.1\%{]} \\
\end{longtable}
}

\emph{Table 2: Per-category FP rates, val split (BCE baseline, seed 42,
τ=0.5). River is the only category with confirmed non-zero FP across
both models. EfficientNetB0 has higher river FP (9.21\% vs.~3.95\%) but
better flood recall at τ=0.5 (97.8\% vs.~97.2\%). Bold CI upper bounds
for swimming pool (12.3\%) and plant (10.3\%) indicate these categories
require \textasciitilde150 and \textasciitilde120 val images,
respectively, for a 7\%-wide 95\% CI.}

\textbf{Confounder heatmap.} Figure 5 shows a heatmap of per-category FP
rates across all evaluated conditions. Only the river column contains
non-zero cells; all other categories are zero across all models and
ablation conditions.

\begin{figure}
\centering
\pandocbounded{\includegraphics[keepaspectratio,alt={Per-category confounder FP heatmap}]{../results/figures/confounder_fp_heatmap.pdf}}
\caption{Per-category confounder FP heatmap}
\end{figure}

\emph{Figure 5: Per-category false positive rate heatmap across all
conditions (val set, seed 42). River is the only confounder category
requiring targeted remediation. Swimming pool is statistically
underpowered (N=28; 0\% FP does not imply robustness).}

\begin{quote}
\textbf{Statistical power note on swimming pool}: The observed 0\% FP on
28 swimming pool val images has an upper 95\% CI bound of
\textbf{12.3\%}. We cannot declare either model robust to swimming pool
confounders from this data. Approximately 150 val images would be
required to reduce the 95\% CI upper bound below 2.5\% at 0\% observed
FP.
\end{quote}

\textbf{The river result directly motivates Phase-1 HNM targeting.}
EfficientNetB0's 9.21\% river FP rate (95\% CI: {[}3.8\%, 17.7\%{]}) is
high enough to flag as a mining target (above the θ=0.05 threshold).
ResNet50's 3.95\% rate is also above threshold but EfficientNetB0 is the
primary HNM target given its higher baseline river FP.

\subsection{5.6 Phase-1 Hard Negative Mining
Results}\label{phase-1-hard-negative-mining-results}

Table 3 reports all HNM ablation results for EfficientNetB0 BCE (seed
42). Phase-1 HNM achieves the best result on both primary objectives
simultaneously.

{\def\LTcaptype{none} % do not increment counter
\begin{longtable}[]{@{}
  >{\raggedright\arraybackslash}p{(\linewidth - 12\tabcolsep) * \real{0.1429}}
  >{\raggedright\arraybackslash}p{(\linewidth - 12\tabcolsep) * \real{0.1429}}
  >{\raggedright\arraybackslash}p{(\linewidth - 12\tabcolsep) * \real{0.1429}}
  >{\raggedright\arraybackslash}p{(\linewidth - 12\tabcolsep) * \real{0.1429}}
  >{\raggedright\arraybackslash}p{(\linewidth - 12\tabcolsep) * \real{0.1429}}
  >{\raggedright\arraybackslash}p{(\linewidth - 12\tabcolsep) * \real{0.1429}}
  >{\raggedright\arraybackslash}p{(\linewidth - 12\tabcolsep) * \real{0.1429}}@{}}
\toprule\noalign{}
\begin{minipage}[b]{\linewidth}\raggedright
Condition
\end{minipage} & \begin{minipage}[b]{\linewidth}\raggedright
Val Acc ↑
\end{minipage} & \begin{minipage}[b]{\linewidth}\raggedright
Flood Recall ↑
\end{minipage} & \begin{minipage}[b]{\linewidth}\raggedright
FN ↓
\end{minipage} & \begin{minipage}[b]{\linewidth}\raggedright
PR-AUC
\end{minipage} & \begin{minipage}[b]{\linewidth}\raggedright
River FP Rate
\end{minipage} & \begin{minipage}[b]{\linewidth}\raggedright
McNemar p
\end{minipage} \\
\midrule\noalign{}
\endhead
\bottomrule\noalign{}
\endlastfoot
BCE baseline & 98.29\% & 97.8\% & 7 (1 in 46) & 0.9976 & 9.21\% (7/76) &
--- \\
Random injection & 98.41\% & 98.1\% & 6 (1 in 54) & 0.9981 & 9.21\%
(7/76) & p \textgreater{} 0.05 \\
No injection (extended) & 98.66\% & 98.1\% & 6 (1 in 54) & 0.9985 &
6.58\% (5/76) & p \textgreater{} 0.05 \\
\textbf{Phase-1 HNM-BCE ★} & \textbf{98.78\%} & \textbf{99.1\%} &
\textbf{3 (1 in 107)} & \textbf{0.9981} & \textbf{5.26\% (4/76)} &
\textbf{p \textgreater{} 0.05} \\
\end{longtable}
}

\emph{Table 3: EfficientNetB0 BCE HNM ablation results (seed 42 only;
multi-seed bootstrap CI pending). Phase-1 HNM achieves best overall
accuracy (98.78\%) and best flood recall (99.1\%, 3 FN), while also
reducing river FP from 9.21\%→5.26\% (7→4 FP). The no-injection control
improves river FP (6.58\%) without improving flood recall beyond random
injection. Only HNM simultaneously reduces both FN (7→3) and river FP
(7→4). McNemar's test (Bonferroni-corrected) shows p \textgreater{} 0.05
for all EfficientNetB0 pairs --- underpowered at N=76 river val images
and N=322 flood val images.}

\textbf{Key result: Phase-1 HNM uniquely improves both objectives.}
Random injection (category exposure without difficulty ranking) reduces
FN from 7 to 6 but does not reduce river FP. Extended training (no
injection) reduces river FP from 9.21\% to 6.58\% but also leaves FN at
6. Only Phase-1 HNM achieves both: FN 7→3 (57\% reduction) and river FP
9.21\%→5.26\% (43\% reduction). After HNM, EfficientNetB0's river FP
rate (5.26\%, 4/76) approaches ResNet50's baseline rate (3.95\%, 3/76)
while substantially improving flood recall (99.1\% vs.~ResNet50's
97.2\%).

\textbf{Difficulty ranking is the key lever.} The comparison between
random injection (6 FN, 9.21\% river FP) and HNM (3 FN, 5.26\% river FP)
with the same number of injected images isolates the contribution of
probability-ranked candidate selection. Both inject images from the
river category; only HNM selects the hardest-scoring ones.

\textbf{ResNet50 HNM ablations pending re-evaluation.} ResNet50 HNM
experiments previously conducted under a different baseline
configuration are excluded from this table. ResNet50 ablations
(no-injection, random injection, HNM-BCE, HNM-Focal) will be
re-evaluated against the corrected ResNet50 BCE baseline reported in
Table 1.

\subsection{5.7 Calibration Analysis}\label{calibration-analysis}

\emph{Pending. Temperature scaling and reliability diagrams to be added
after HNM runs complete.}

{\def\LTcaptype{none} % do not increment counter
\begin{longtable}[]{@{}
  >{\raggedright\arraybackslash}p{(\linewidth - 8\tabcolsep) * \real{0.2000}}
  >{\raggedright\arraybackslash}p{(\linewidth - 8\tabcolsep) * \real{0.2000}}
  >{\raggedright\arraybackslash}p{(\linewidth - 8\tabcolsep) * \real{0.2000}}
  >{\raggedright\arraybackslash}p{(\linewidth - 8\tabcolsep) * \real{0.2000}}
  >{\raggedright\arraybackslash}p{(\linewidth - 8\tabcolsep) * \real{0.2000}}@{}}
\toprule\noalign{}
\begin{minipage}[b]{\linewidth}\raggedright
Model
\end{minipage} & \begin{minipage}[b]{\linewidth}\raggedright
Condition
\end{minipage} & \begin{minipage}[b]{\linewidth}\raggedright
ECE (pre-calibration) ↓
\end{minipage} & \begin{minipage}[b]{\linewidth}\raggedright
ECE (post-calibration) ↓
\end{minipage} & \begin{minipage}[b]{\linewidth}\raggedright
Temperature T
\end{minipage} \\
\midrule\noalign{}
\endhead
\bottomrule\noalign{}
\endlastfoot
EfficientNetB0 & BCE baseline & --- & --- & --- \\
EfficientNetB0 & Phase-1 HNM & --- & --- & --- \\
ResNet50 & BCE baseline & --- & --- & --- \\
\end{longtable}
}

\emph{Table 4: Calibration results (pending). ECE = Expected Calibration
Error; lower is better.}

\subsection{5.8 Multi-Seed Results}\label{multi-seed-results}

\emph{Pending. Seeds 1, 7, 123, 2024 in addition to seed 42.}

{\def\LTcaptype{none} % do not increment counter
\begin{longtable}[]{@{}
  >{\raggedright\arraybackslash}p{(\linewidth - 8\tabcolsep) * \real{0.2000}}
  >{\raggedright\arraybackslash}p{(\linewidth - 8\tabcolsep) * \real{0.2000}}
  >{\raggedright\arraybackslash}p{(\linewidth - 8\tabcolsep) * \real{0.2000}}
  >{\raggedright\arraybackslash}p{(\linewidth - 8\tabcolsep) * \real{0.2000}}
  >{\raggedright\arraybackslash}p{(\linewidth - 8\tabcolsep) * \real{0.2000}}@{}}
\toprule\noalign{}
\begin{minipage}[b]{\linewidth}\raggedright
Model
\end{minipage} & \begin{minipage}[b]{\linewidth}\raggedright
Condition
\end{minipage} & \begin{minipage}[b]{\linewidth}\raggedright
Flood Recall (mean ± 95\% CI)
\end{minipage} & \begin{minipage}[b]{\linewidth}\raggedright
PR-AUC (mean ± 95\% CI)
\end{minipage} & \begin{minipage}[b]{\linewidth}\raggedright
River FP Rate (mean ± 95\% CI)
\end{minipage} \\
\midrule\noalign{}
\endhead
\bottomrule\noalign{}
\endlastfoot
EfficientNetB0 & BCE baseline & --- & --- & --- \\
EfficientNetB0 & Phase-1 HNM & --- & --- & --- \\
ResNet50 & BCE baseline & --- & --- & --- \\
\end{longtable}
}

\emph{Table 5: Multi-seed results (n=5 seeds, 95\% bootstrap CI).
Pending.}

\begin{center}\rule{0.5\linewidth}{0.5pt}\end{center}

\section{6. Discussion}\label{discussion}

\subsection{6.1 Both Architectures Achieve Strong Baselines --- With
Different
Tradeoffs}\label{both-architectures-achieve-strong-baselines-with-different-tradeoffs}

The corrected ResNet50 BCE baseline (PR-AUC=0.9991, 97.2\% flood recall,
9 FN) challenges the naive assumption that a larger architecture is
always worse on small datasets. EfficientNetB0 (PR-AUC=0.9976, 97.8\%
flood recall, 7 FN) achieves better per-threshold flood recall at τ=0.5
while incurring more river false positives (9.21\% vs.~3.95\%). The
architectures present a genuine operating point tradeoff: ResNet50's
lower FP makes it preferable for false-alarm-sensitive deployment
contexts; EfficientNetB0's higher recall makes it preferable for
recall-sensitive screening.

This tradeoff underscores why per-threshold performance and per-category
FP analysis are both required for deployment decisions --- PR-AUC alone
does not reveal which operating point each architecture excels at.

\subsection{6.2 Phase-1 HNM Resolves the EfficientNetB0
Tradeoff}\label{phase-1-hnm-resolves-the-efficientnetb0-tradeoff}

Phase-1 HNM applied to EfficientNetB0 BCE (the architecture with higher
recall but higher river FP) simultaneously improves both metrics: recall
97.8\%→99.1\% and river FP 9.21\%→5.26\%. This resolves the architecture
tradeoff: post-HNM EfficientNetB0 achieves 99.1\% flood recall
(vs.~ResNet50's 97.2\%) and 5.26\% river FP (approaching ResNet50's
3.95\%), combining the strengths of both baselines. The 3 remaining
flood misses (1 in 107) represent a substantial improvement from the
7-miss baseline (1 in 46).

Critically, this improvement requires both the right mining signal
(Phase 1, not Phase 2) and the right selection mechanism (difficulty
ranking, not random sampling). Neither random injection nor extended
training achieves the combined improvement --- each only partially
solves the problem.

\subsection{6.3 The Phase 1 Mining Design: Why It
Works}\label{the-phase-1-mining-design-why-it-works}

\textbf{Memorisation destroys mining signal.} By end of Phase 2,
EfficientNetB0 assigns near-zero flood probability to all training
non-flood images: the observed training-set river FP rate is 0.31\%.
Mining from this checkpoint produces no useful candidates --- the model
has already learned to reject all training negatives, even the visually
confusing ones. The Phase 1 checkpoint, with only the classification
head and top 30 backbone layers trained, retains genuine uncertainty on
river and swimming pool images. This residual confusion is the mining
signal that makes candidate selection meaningful.

\textbf{Phase 1 mining is not simply using a weaker model.} The key
property is that Phase 1 uncertainty is \emph{correlated with visual
confound difficulty}: the images the Phase 1 model is most uncertain
about (high p̂ despite being non-flood) are genuinely harder to classify.
Random injection (using the same images but without difficulty ranking)
produces inferior results, confirming that the ranked difficulty signal,
not just category exposure, is the mechanism of improvement.

\subsection{6.4 PR-AUC vs.~ROC-AUC for
Screening}\label{pr-auc-vs.-roc-auc-for-screening}

Both architectures achieve high ROC-AUC (EfficientNetB0: 0.9985,
ResNet50: 0.9995) and high PR-AUC (0.9976 and 0.9991, respectively). In
the corrected baseline setting, the gap between architectures is small
on both metrics. This scenario --- where both metrics agree that both
models are strong --- is the expected outcome when both models are
properly configured. The key diagnostic value of PR-AUC emerges when a
model has systematic recall collapse; in our setting, it primarily
disambiguates the architecture tradeoff at high recall operating points,
where EfficientNetB0's curve shows better-maintained precision above
recall ≈ 0.97.

\subsection{6.5 Focal Loss and
Architecture}\label{focal-loss-and-architecture}

EfficientNetB0 Focal (98.5\% recall, 5 FN) outperforms EfficientNetB0
BCE (97.8\% recall, 7 FN), suggesting that EfficientNetB0's
well-calibrated feature space benefits from within-batch hard-example
reweighting. ResNet50 Focal results pending re-evaluation; the
previously observed degradation under the uncorrected baseline may not
replicate with the corrected preprocessing configuration.

\subsection{6.6 Limitations}\label{limitations}

\begin{itemize}
\tightlist
\item
  \textbf{Small dataset}: 4,099 images; per-category val counts as low
  as 28 (swimming pool). Clopper-Pearson 95\% CI upper bounds highlight
  the limits of inference at small N.
\item
  \textbf{Underpowered confounder comparisons}: With N=76 river val
  images, McNemar's test has \textasciitilde30\% power to detect a 3--4
  percentage point FP change. Approximately 300 river val images would
  be required for 80\% power.
\item
  \textbf{Single geographic/source distribution}: Dataset provenance
  (US-centric) may not generalise to other geographic regions or urban
  morphologies.
\item
  \textbf{Phase boundary not ablated}: The (30, 50) freeze boundary is a
  design choice; optimal values for flood data are not evaluated.
\item
  \textbf{No external test set}: Generalisation to AlleyFloodNet and
  other out-of-distribution benchmarks is not evaluated.
\item
  \textbf{Multi-seed results pending}: Seed 42 results should be
  interpreted as single-run estimates; bootstrap CIs across 5 seeds will
  bound seed sensitivity.
\end{itemize}

\begin{center}\rule{0.5\linewidth}{0.5pt}\end{center}

\section{7. Conclusion}\label{conclusion}

We presented a Phase-1 Hard Negative Mining protocol for crowdsourced
flood image screening that addresses the dual challenge of maximising
flood recall and minimising per-category confounder false positive
rates. Our benchmark of 4,099 images across 8 confounder categories
(after SHA-256 deduplication) establishes that --- when preprocessing is
correctly configured --- both EfficientNetB0 (PR-AUC 0.9976, 97.8\%
flood recall) and ResNet50 (PR-AUC 0.9991, 97.2\% flood recall) achieve
strong baselines, with the two architectures presenting a meaningful
recall vs.~precision tradeoff on river confounders.

Phase-1 HNM resolves this tradeoff for EfficientNetB0: by mining
difficulty-ranked hard negatives from the partially-trained Phase 1
checkpoint --- where the model retains genuine confusion on river images
--- and injecting them into retraining from the converged Phase 2
baseline, we simultaneously improve flood recall from 97.8\% to 99.1\%
(FN 7→3, 1 in 46 → 1 in 107) and reduce river false positives from
9.21\% to 5.26\% (FP 7→4). Extended training without injection and
random injection with the same image count both fail to achieve this
joint improvement, confirming that difficulty-ranked candidate selection
is the operative mechanism.

Three practical recommendations emerge from this work. First,
\textbf{per-category false positive rates} should be reported alongside
aggregate metrics in flood screening evaluations; river confounders are
the dominant failure mode and are invisible in aggregate accuracy or
PR-AUC reports. Second, \textbf{Phase-1 mining} should be preferred over
Phase-2 mining for training-distribution-augmentation tasks: the
intermediate checkpoint retains meaningful uncertainty that the
converged model has memorised away. Third, \textbf{EfficientNetB0 +
Phase-1 HNM} is our recommended screening configuration, achieving
99.1\% flood recall with river FP reduced to 5.26\%, suitable for
deployment as an automated pre-filter before human analyst review.

Multi-seed bootstrap confidence intervals, ResNet50 HNM ablations,
probability calibration, and external test set evaluation
(AlleyFloodNet) are pending and will complete the experimental picture.

\begin{center}\rule{0.5\linewidth}{0.5pt}\end{center}

\section{References}\label{references}

Alam, F., Alam, T., Hasan, M. A., Hasnat, A., Imran, M., \& Ofli, F.
(2023). MEDIC: A multi-task learning dataset for disaster image
classification. \emph{Neural Computing and Applications}, \emph{35}(3),
2609--2632. https://doi.org/10.1007/s00521-022-07717-0

Alam, F., Imran, M., \& Ofli, F. (2018). CrisisMMD: Multimodal Twitter
datasets from natural disasters. \emph{Proceedings of ICWSM 2018},
465--473. https://doi.org/10.1609/icwsm.v12i1.14983

Alam, F., Ofli, F., Imran, M., Alam, T., \& Qazi, U. (2021).
CrisisBench: Benchmarking crisis-related social media datasets for
humanitarian information processing. \emph{Proceedings of ICWSM 2021},
923--934. https://doi.org/10.1609/icwsm.v15i1.18115

Cai, G. (2025). Prioritizing recall: Recall-first machine learning for
traffic accident severity detection. \emph{IRJAEM}, \emph{3}(6).
https://doi.org/10.61173/ym16tb60

Clopper, C. J., \& Pearson, E. S. (1934). The use of confidence or
fiducial limits illustrated in the case of the binomial.
\emph{Biometrika}, \emph{26}(4), 404--413.
https://doi.org/10.1093/biomet/26.4.404

Davis, J., \& Goadrich, M. (2006). The relationship between
precision-recall and ROC curves. \emph{Proceedings of ICML 2006},
233--240. https://doi.org/10.1145/1143844.1143874

Dong, L., Zhao, W., Lan, Z., Liu, Y., \& Shi, X. (2025). Large-small
model synergy for high-recall UAV-based aerial surveillance.
\emph{Proceedings of AIHCIR 2025}.
https://doi.org/10.1109/AIHCIR67580.2025.11404852

Dong, J., Jiang, Z., Pan, D., Chen, Z., Guan, Q., Zhang, H., Gui, G., \&
Gui, W. (2025). A survey on confidence calibration of deep
learning-based classification models under class imbalance data.
\emph{IEEE Transactions on Neural Networks and Learning Systems}.
https://doi.org/10.1109/TNNLS.2025.3565159

Dosovitskiy, A., Beyer, L., Kolesnikov, A., Weissenborn, D., Zhai, X.,
Unterthiner, T., Dehghani, M., Minderer, M., Heigold, G., Gelly, S.,
Uszkoreit, J., \& Houlsby, N. (2021). An image is worth 16x16 words:
Transformers for image recognition at scale. \emph{Proceedings of ICLR
2021}. https://arxiv.org/abs/2010.11929

Efron, B., \& Tibshirani, R. J. (1993). \emph{An Introduction to the
Bootstrap}. Chapman \& Hall.

Esparza, M., Farahmand, H., Brody, S., \& Mostafavi, A. (2022).
Examining data imbalance in crowdsourced reports for improving flash
flood situational awareness. \emph{ArXiv preprint arXiv:2207.05797}.
https://arxiv.org/abs/2207.05797

Guo, C., Pleiss, G., Sun, Y., \& Weinberger, K. Q. (2017). On
calibration of modern neural networks. \emph{Proceedings of ICML 2017},
1321--1330. https://dl.acm.org/doi/10.5555/3305381.3305518

He, K., Zhang, X., Ren, S., \& Sun J. (2016). Deep residual learning for
image recognition. \emph{Proceedings of CVPR 2016}, 770--778.
https://doi.org/10.1109/CVPR.2016.90

Khan, M. A., Ahmed, N., Padela, J., Raza, M. S., Gangopadhyay, A., Wang,
J., Foulds, J., Busart, C. E., \& Erbacher, R. (2023). Flood-ResNet50:
Optimized deep learning model for efficient flood detection on edge
device. \emph{Proceedings of ICMLA 2023}, 506--511.
https://doi.org/10.1109/ICMLA58977.2023.00077

Kirkpatrick, J., Pascanu, R., Rabinowitz, N., Veness, J., Desjardins,
G., Rusu, A. A., Milan, K., Quan, J., Ramalho, T., Grabska-Barwinska,
A., Hassabis, D., Clopath, C., Kumaran, D., \& Hadsell, R. (2017).
Overcoming catastrophic forgetting in neural networks. \emph{Proceedings
of the National Academy of Sciences}, \emph{114}(13), 3521--3526.
https://doi.org/10.1073/pnas.1611835114

Lee, S., Park, J., Kim, D., \& others (2025). AlleyFloodNet: A
ground-level image dataset for rapid flood detection in economically and
flood-vulnerable areas. \emph{Electronics}, \emph{14}(10), 2082.
https://doi.org/10.3390/electronics14102082

Lin, T.-Y., Goyal, P., Girshick, R., He, K., \& Dollár, P. (2017). Focal
loss for dense object detection. \emph{Proceedings of ICCV 2017},
2980--2988. https://doi.org/10.1109/ICCV.2017.324

Lyu, S., Lyu, R., Zhao, Y., \& Gao, W. (2025). ResNet50-based
progressive transfer learning. \emph{Proceedings of ICCNSE 2025}.
https://doi.org/10.1109/ICCNSE66404.2025.11144108

Neupane, B., Aryal, J., \& Rajabifard, A. (2025). Fine-tuning-based
transfer learning for building extraction. \emph{Remote Sensing},
\emph{17}(7), 1251. https://doi.org/10.3390/rs17071251

Papadimos, T., Andreadis, S., Gialampoukidis, I., Vrochidis, S., \&
Kompatsiaris, I. (2023). Flood-related multimedia benchmark evaluation:
Challenges, results and a novel GNN approach. \emph{Applied Sciences},
\emph{13}(8), 4814. https://doi.org/10.3390/app13084814

Rahnemoonfar, M., Chowdhury, T., Sarkar, A., Varshney, D., Yari, M., \&
Murphy, R. (2020). FloodNet: A high resolution aerial imagery dataset
for post flood scene understanding. \emph{ArXiv preprint
arXiv:2012.02951}. https://arxiv.org/abs/2012.02951

Saito, T., \& Rehmsmeier, M. (2015). The precision-recall plot is more
informative than the ROC plot when evaluating binary classifiers on
imbalanced datasets. \emph{PLOS ONE}, \emph{10}(3), e0118432.
https://doi.org/10.1371/journal.pone.0118432

Shrivastava, A., Gupta, A., \& Girshick, R. (2016). Training
region-based object detectors with online hard example mining.
\emph{Proceedings of CVPR 2016}, 761--769.
https://arxiv.org/abs/1604.03540

Tan, M., \& Le, Q. V. (2019). EfficientNet: Rethinking model scaling for
convolutional neural networks. \emph{Proceedings of ICML 2019},
6105--6114. https://proceedings.mlr.press/v97/tan19a.html

Tellman, B., Sullivan, J. A., Kuhn, C., Kettner, A. J., Doyle, C. S.,
Brakenridge, G. R., Erickson, T. A., \& Slayback, D. A. (2021).
Satellite imaging reveals increased proportion of population exposed to
floods. \emph{Nature}, \emph{596}, 80--86.
https://doi.org/10.1038/s41586-021-03695-w

Viola, P., \& Jones, M. (2004). Robust real-time face detection.
\emph{International Journal of Computer Vision}, \emph{57}(2), 137--154.
https://doi.org/10.1023/B:VISI.0000013087.49260.fb

\begin{center}\rule{0.5\linewidth}{0.5pt}\end{center}

\section{Appendix}\label{appendix}

\subsection{A. Hyperparameters}\label{a.-hyperparameters}

{\def\LTcaptype{none} % do not increment counter
\begin{longtable}[]{@{}ll@{}}
\toprule\noalign{}
Hyperparameter & Value \\
\midrule\noalign{}
\endhead
\bottomrule\noalign{}
\endlastfoot
Phase 1 LR & 1×10⁻³ \\
Phase 2 LR & 1×10⁻⁴ \\
Phase 1 max epochs & 30 \\
Phase 2 max epochs & 50 \\
Early stopping patience & 10 (val loss) \\
LR reduce patience (Phase 1) & 5 \\
LR reduce patience (Phase 2) & 3 \\
Batch size & 32 \\
Image size & 224×224 \\
Focal loss γ & 2.0 \\
HNM mining checkpoint & Phase 1 \\
HNM retraining start & Phase 2 checkpoint \\
HNM top\_pct & 0.20 \\
HNM FP threshold θ & 0.05 \\
HNM augmentation multiplier & 5× \\
Seeds & 42, 1, 7, 123, 2024 \\
Compute & Apple M4 (16GB, Metal GPU) \\
\end{longtable}
}

\subsection{B. Clopper-Pearson CI Reference
Table}\label{b.-clopper-pearson-ci-reference-table}

Upper bound of 95\% CI for observed FP rate = 0\%, by sample size.

{\def\LTcaptype{none} % do not increment counter
\begin{longtable}[]{@{}lll@{}}
\toprule\noalign{}
N val images & Upper 95\% CI bound & Category \\
\midrule\noalign{}
\endhead
\bottomrule\noalign{}
\endlastfoot
10 & 30.8\% & --- \\
15 & 21.8\% & --- \\
28 & \textbf{12.3\%} & swimming pool \\
34 & \textbf{10.3\%} & plant \\
50 & 7.1\% & vehicle \\
55 & 6.5\% & building \\
67 & 5.4\% & animal \\
76 & 4.7\% & river (if observed = 0) \\
82 & 4.4\% & park\_walkway \\
105 & 3.5\% & street\_clear \\
\textasciitilde150 & \textasciitilde2.0\% & needed for swimming pool
meaningful bound \\
\textasciitilde300 & \textasciitilde1.0\% & needed for river meaningful
bound \\
\end{longtable}
}

\subsection{C. Dataset Release}\label{c.-dataset-release}

Training dataset: \texttt{zinnia82/flood-binary-hnm-benchmark} (Hugging
Face Hub). Code: {[}TODO --- GitHub link before submission{]}.

\subsection{D. What Changed from Earlier
Draft}\label{d.-what-changed-from-earlier-draft}

The following changes from the previous draft (paper.md) reflect updated
experimental results and corrections:

{\def\LTcaptype{none} % do not increment counter
\begin{longtable}[]{@{}
  >{\raggedright\arraybackslash}p{(\linewidth - 4\tabcolsep) * \real{0.3333}}
  >{\raggedright\arraybackslash}p{(\linewidth - 4\tabcolsep) * \real{0.3333}}
  >{\raggedright\arraybackslash}p{(\linewidth - 4\tabcolsep) * \real{0.3333}}@{}}
\toprule\noalign{}
\begin{minipage}[b]{\linewidth}\raggedright
Section
\end{minipage} & \begin{minipage}[b]{\linewidth}\raggedright
Change
\end{minipage} & \begin{minipage}[b]{\linewidth}\raggedright
Reason
\end{minipage} \\
\midrule\noalign{}
\endhead
\bottomrule\noalign{}
\endlastfoot
Abstract & ResNet50 results corrected to PR-AUC=0.9991, 97.2\% recall (9
FN) & Preprocessing fix \\
§5.2 Table 1 & ResNet50 BCE: was (90.2\% acc, 80.4\% recall, 63 FN); now
(98.5\% acc, 97.2\% recall, 9 FN) & Preprocessing fix \\
§5.6 Table 3 & EfficientNetB0 HNM river FP updated to 5.26\% (4/76); was
9.21\% placeholder in draft & New confounder CSV
(\texttt{confounder\_fp\_rates\_efficientnet\_p1hnm\_VAL.csv}) \\
§6 Discussion & Architecture comparison reframed from ``ResNet50 fails''
to ``genuine tradeoff'' & Corrected baseline \\
§7 Conclusion & Updated to reflect joint improvement from HNM (both
recall and FP) & Updated Table 3 \\
\end{longtable}
}

\end{document}
